\subsection{Requested-temperature sensitivity}
\label{app:hosted-temperature}
Grok 4.3 and Kimi K3 each receive requested temperatures $T=1.0$ and
$T=1.5$ on the same 120 tasks: eight prompts in each of five domains and
three levels. Eight responses from separate requests per prompt give 960
responses per deployment--temperature condition and 3,840 in total.
Within each deployment, the two conditions share all non-temperature
request settings and run concurrently. Accuracy and mode counts include
all responses, including verification failures and truncated outputs.

The providers accept both requested temperature settings, but their responses
do not report effective sampling temperatures. DeepSeek is outside this
comparison because its thinking-mode interface ignores the temperature setting.
The results concern these requested settings on this 120-task subset and
do not establish a general temperature response for all hosted deployments.

Each level summary averages five domains equally; All averages the 15
domain--level cells equally. Collision pools correct pairs within each cell,
so tasks with more correct responses receive more weight. Eligible tasks
and pair weights can differ between temperatures; the contrast does not
hold accuracy or the correct-pair population fixed. An undefined cell makes
the corresponding average undefined. Pointwise 95\% intervals use 20,000
paired whole-prompt bootstrap resamples, stratified by domain and level.
Each resampled prompt includes all eight responses. Prompts are paired
across conditions; output draws are not. Intervals have no multiplicity adjustment.

\paragraph{Observed sensitivity.}
For Grok, the observed diversity change is +0.050 [$-$0.058, +0.158] modes, and the collision change is $-$0.78 [$-$3.86, +2.12] percentage points. Both intervals include zero, so neither shows a clear diversity gain.
The per-response accuracy change is $-$2.29 [$-$3.85, $-$0.94] percentage points. These results do not establish invariance to temperature in general.

At $T=1.5$, 696 of Kimi's 960 outputs reach the token limit,
compared with 0 at $T=1.0$. Strict accuracy falls from 96.56\% to 21.25\%, while raw diversity falls from 1.933 to 0.983 modes.
The large increase in truncation accompanies losses in accuracy and distinct modes.
These results do not isolate correct-output concentration at matched accuracy.
The five-domain collision mean is undefined at
$T=1.5$ because some cells have no eligible correct pairs. Neither model
has provider-declared refusals in this experiment.

$\dagger$ marks a defined contrast whose interval is omitted because some
bootstrap resamples have undefined cells. Five-domain averages require all
five cell estimates; conditioning on only the estimable resamples would
change the uncertainty calculation.

\begin{table}[!htbp]
  \centering
  \setlength{\parfillskip}{0pt plus .20\linewidth}
  \caption{\textbf{Higher requested temperature yields no clear Grok diversity gain and much lower Kimi accuracy.} Strict executable grades compare $T=1.5$ with $T=1.0$.
  Each condition contains eight responses on each of 120 tasks.
  Level summaries average five domains equally; All averages 15 domain--level cells.
  A is per-response accuracy (\%); D is mean \texttt{distinct@8};
  C is correct-pair collision (\%), pooling pairs within each cell.
  Differences subtract $T=1.0$ from $T=1.5$: mode counts for D and percentage points for A/C.
  Brackets give pointwise 95\% intervals from paired whole-prompt bootstrap resampling.
  $\dagger$ marks an omitted interval when some resamples have undefined cells;
  -- denotes an undefined estimate or omitted interval.}
  \small
  \setlength{\tabcolsep}{5pt}
  \begin{tabular}{@{}llrrrr@{}}
    \toprule
    Deployment & Level & Metric & $T=1.0$ & $T=1.5$ & Difference [95\% interval] \\
    \midrule
    Grok 4.3 & 1 & A & 98.12 & 95.62 & $-$2.50 [$-$4.06, $-$0.94] \\
     &  & D & 1.900 & 2.000 & +0.100 [$-$0.100, +0.300] \\
     &  & C & 72.89 & 73.16 & +0.27 [$-$6.31, +6.37] \\
     & 2 & A & 99.06 & 99.38 & +0.31 [$-$0.94, +1.56] \\
     &  & D & 1.425 & 1.425 & +0.000 [$-$0.175, +0.200] \\
     &  & C & 86.61 & 85.20 & $-$1.41 [$-$6.26, +3.30] \\
     & 3 & A & 98.75 & 94.06 & $-$4.69 [$-$9.06, $-$0.94] \\
     &  & D & 1.500 & 1.550 & +0.050 [$-$0.125, +0.225] \\
     &  & C & 84.04 & 82.85 & $-$1.20 [$-$5.48, +2.79] \\
     & All & A & 98.65 & 96.35 & $-$2.29 [$-$3.85, $-$0.94] \\
     &  & D & 1.608 & 1.658 & +0.050 [$-$0.058, +0.158] \\
     &  & C & 81.18 & 80.40 & $-$0.78 [$-$3.86, +2.12] \\
    \midrule
    Kimi K3 & 1 & A & 92.50 & 22.50 & $-$70.00 [$-$75.62, $-$63.75] \\
     &  & D & 2.200 & 1.075 & $-$1.125 [$-$1.475, $-$0.775] \\
     &  & C & 66.23 & 53.12 & $-$13.11 [--]$^\dagger$ \\
     & 2 & A & 98.12 & 20.31 & $-$77.81 [$-$82.50, $-$72.81] \\
     &  & D & 1.850 & 0.900 & $-$0.950 [$-$1.300, $-$0.625] \\
     &  & C & 72.18 & 66.55 & $-$5.63 [--]$^\dagger$ \\
     & 3 & A & 99.06 & 20.94 & $-$78.12 [$-$83.75, $-$72.81] \\
     &  & D & 1.750 & 0.975 & $-$0.775 [$-$1.025, $-$0.525] \\
     &  & C & 78.33 & -- & -- \\
     & All & A & 96.56 & 21.25 & $-$75.31 [$-$78.44, $-$72.19] \\
     &  & D & 1.933 & 0.983 & $-$0.950 [$-$1.133, $-$0.767] \\
     &  & C & 72.25 & -- & -- \\
    \bottomrule
  \end{tabular}
\end{table}

\begin{table}[!htbp]
  \centering
  \setlength{\parfillskip}{0pt plus .20\linewidth}
  \caption{\textbf{Formatting normalization leaves the temperature sensitivity largely unchanged.} Only rows differing from strict grading appear; all other rows are identical.
  Each condition contains eight responses on each of 120 tasks.
  Level summaries average five domains equally; All averages 15 domain--level cells.
  A is per-response accuracy (\%); D is mean \texttt{distinct@8};
  C is correct-pair collision (\%), pooling pairs within each cell.
  Differences subtract $T=1.0$ from $T=1.5$: mode counts for D and percentage points for A/C.
  Brackets give pointwise 95\% intervals from paired whole-prompt bootstrap resampling.
  $\dagger$ marks an omitted interval when some resamples have undefined cells;
  -- denotes an undefined estimate or omitted interval.}
  \small
  \setlength{\tabcolsep}{5pt}
  \begin{tabular}{@{}llrrrr@{}}
    \toprule
    Deployment & Level & Metric & $T=1.0$ & $T=1.5$ & Difference [95\% interval] \\
    \midrule
    Kimi K3 & 1 & A & 92.50 & 22.81 & $-$69.69 [$-$75.31, $-$63.44] \\
     & 1 & C & 66.23 & 53.61 & $-$12.62 [--]$^\dagger$ \\
     & 3 & A & 99.38 & 20.94 & $-$78.44 [$-$84.06, $-$73.12] \\
     & 3 & C & 78.18 & -- & -- \\
     & All & A & 96.67 & 21.35 & $-$75.31 [$-$78.44, $-$72.19] \\
     & All & C & 72.20 & -- & -- \\
    \bottomrule
  \end{tabular}
\end{table}

Across the hosted comparisons, \texttt{Countdown} has no uniform reference because its
total support is not certified, so uniform-reference comparisons use the other
four domains. Cross-level differences
also reflect different task populations and do not establish a common ordering
of model difficulty. These inference-only results do not identify a training
or parameter-scale cause of concentration.
